\head{第14节后半：有限点、平展性、双有理投影及普遍图式}
\source{法文手稿物理88页第二段至90页第15节标题之前；英文物理57页More delicate起至59页第15节标题之前；旧注V-H-L029--047}

\audit{V22-DIF-E01}{原稿、手改与英文转写的归属}{三页法稿及对应英稿已实际看图，正文不再沿用可恢复处的“字迹不清”}
法88的微分性质导言仍有“表明／若未作则将会表明”的草稿语气；本版没有把普通奇点、唯一性和有理性整组断言写成已重新证明。法稿将证明留给Dieudonné，英文括加“或Blass”属编者插语。14.4(a)中“单态射”已经被作者划去并改为“根态射”；该删改必须保留其方向。

法88末行14.4(b)的“甚至在整条逆像的各点”已经可辨；法89的14.5平展性补句也可辨。二者文字恢复后仍须数学检查，不能把它们从“未辨读”直接改成“成立”。下面分别给出反例。14.4(a)的空交式确实没有附加不交条件；14.5的闭点也确实没有光滑限定，均不能归罪于英文。

法89可恢复的读法还有：$p'_{C_\eta}$的撇号、“譬如假设$k$代数闭”、以及“或者，在维数$m$时还有更明显的反例”。因此旧注中疑似$m$与$m-1$矛盾的问题由完整连接语消除。法90首图左下为$S$、第二图右上等号两边都有波浪号，普遍底和普遍中心有下划线；英文漏记这些区别。本版分别恢复并以$\mathcal G,\mathcal C$消除试验底$T$与普遍底的冲突。法90在第15节标题前的全部图式和记号建议均保留。

\audit{V22-DIF-E02}{14.4(a)的两个缺失条件}{零维目标反例，以及不交条件的充分证明}
以下为译者独立核查。取$\mathcal G$为$\mathbb P^r$中余维$m+1$的线性中心的Grassmann概形，$N=\dim\mathcal G$。空中心$m=r$时投影为恒等，下面涉及割线的论证不必使用；先设$0\leq m<r$。

\lead{$m=0$不能省略。}
令$k$为任意域，取$X=T$为$\mathbb P^1_k$中两个不同的$k$有理点，并令$m=0$。$X\hookrightarrow\mathbb P^1$是闭浸入，$\dim X=0$。泛中心是避开这两个点的一个点，投影目标为$\mathbb P^0_{k(\eta)}$，所以两点被合并。这否定原14.4(a)的无条件根态射结论。若$m=0$，$T\to\Spec k$为根态射当且仅当它至多有一个几何点；非约化结构本身不妨碍这一条件。正文采用$m\geq1$的统一充分版本，未将其强加到14.4(b)。

\lead{固定有限点的分离。}
在代数闭包上，将$T$的支集写成有限个几何点。浸入$f$使不同点具有不同的像。对任意两个不同像点$a,b$，在中心避开二者的开集中，投影把它们合并当且仅当中心与直线$\overline{ab}$相交。
设$h=r-m$为中心对应向量子空间的维数。含一个给定向量直线的$h$维子空间，在$\operatorname{Gr}(h,r+1)$中具有余维$m+1$；让该点沿射影直线移动至多增加1维。因此，与一条给定射影直线相交的中心集维数至多$N-m$。当$m\geq1$时它是一个真闭子集。排除有限多个这样的坏闭集，并排除与$X$相交的中心，得到非空开条件。后一个排除成立，因为$\dim X\leq m$，可先用同维的射影闭包。

由此，泛投影在$T$的所有几何点上单射，即其限制为普遍单射。有限非约化$T$也适用，因为判定普遍单射用的是几何点；并未声称它必为闭浸入或单态射。有限点及其几何共轭定义在有限扩域上，坏闭集的有限像仍闭且维数小于$N$，故上述条件能在$k$上表达；论证没有把任意闭点误当成$k$有理点。

\lead{空交式的必要条件与充分性。}
原式只要$Y\cap|T|\ne\varnothing$就失败。例如取$X=\mathbb A^1_k\hookrightarrow\mathbb P^2_k$为直线中的仿射开集，$m=1$，$Y=T=\{0\}$；原维数条件全满足，所列逆像交集却含原点。

现在假设$f$为浸入、$Y\cap|T|=\varnothing$及$\dim Y\leq m-1$。固定一个几何点$t\in T$。考虑$y\in Y$与中心$C$组成的关联集，条件为$C$与割线$\overline{f(t)f(y)}$相交。其维数至多
\[
\dim Y+(N-m)\leq N-1.
\]
即使这个关联集不是固有的，它在$\mathcal G$中的像的闭包仍至多为$N-1$维，因此泛中心不属于该闭包。对有限多个$t$取交，并作上述有限扩域下降，便有
\[
p_{C_\eta}(Y_{k(\eta)})\cap p_{C_\eta}(T_{k(\eta)})=\varnothing.
\]
这与正文空交式等价。$Y$在$X$中闭且$t\notin Y$，故其射影闭包在$X$内仍不含$t$；也可直接对$Y$的有限型关联集计数。不能以“对每个$y$分别一般”为由忽略无穷多个$y$，上式正是统一排除它们所需的维数估计。$m=0$时维数限制迫使$Y$为空，空交式本身没有问题。

\audit{V22-DIF-E03}{14.4(b)的有效部分与手写加强句}{有限指定光滑点的非分歧性成立，整条逆像的加强一般错误}
\lead{指定点处的充分证明。}
在$X$的光滑点$x$，非分歧态射$f$把几何切空间单射到$T_{f(x)}P$。设其维数为$n_x\leq m$。线性投影的微分在这个子空间上单射，当中心避开由该切方向生成的$n_x$维射影切平面时成立。中心的维数为$r-m-1$，所以这样的中心构成非空开集。对有限个几何点同时排除坏条件即可。

这一开条件可下降到$k$：在中心先避开$X$的参数开集上，普遍投影的相对微分层是凝聚层；其支集与$T\times\mathcal G$相交的像由有限态射送到闭集。代数闭包上的切空间计算表明这个闭集是真闭集，故泛中心不在其中。非分歧性可沿域的忠实平坦扩张检验，故结论对非有理点及非约化$T$同样成立。这里光滑性仅在$T$的支集使用。

\lead{整条逆像加强的反例。}
在$\mathbb P^3_k$中取两条异面直线
\[
A=V(X_2,X_3),\qquad L=V(X_0,X_1),
\]
令$D=V(\mathcal I_L^2)$为$L$的一阶无穷小加厚，$X=A\amalg D$取为它们的闭子概形并。二者支集不交，故这个并确是无交并。令$f:X\hookrightarrow\mathbb P^3_k$为闭浸入，$m=1$，并取$T$为$A$上的一个$k$有理点。于是$f$非分歧、$\dim X=1$，且$X$在$T$处光滑，正好满足原14.4(b)的假设。

泛中心是一条同时避开$A,L$的直线。投影在$L$上的限制是
\[
L_{k(\eta)}\xrightarrow{\sim}\mathbb P^1_{k(\eta)},
\]
因为给出投影的两个线性式在$L$上独立；否则有公共零点，中心便与$L$相交。因此$p_{C_\eta}(T)$上必有一个$D$中的点$z$。在代数闭扩域上，$D$在该点的局部模型为
\[
k'[t,u,v]_{(t,u,v)}/(u,v)^2,
\]
其切空间维数为3；目标$\mathbb P^1$的切空间维数为1。微分不可能单射，所以投影在$z$不非分歧。这样，$T$上的正确结论不能加强成整个$p_{C_\eta}^{-1}p_{C_\eta}(T)$上的结论。

这同时解释了为何“分歧集没有固定点”的直观说法必须保留全局光滑的上下文：这里加厚分支提供一个始终存在的坏部分。反例甚至使用闭浸入，故仅给原14.4补上“$f$是浸入”也不能修复这个加强句。

\audit{V22-DIF-E04}{14.5任意闭点处平展的反例}{原稿与英文都有的数学错误，不是记号差错}
取$k$为特征0的代数闭域，令$X\subset\mathbb P^3_k$为仿射曲线
\[
\Spec k[t^3,t^4,t^5]\subset\mathbb A^3_k
\]
的约化射影闭包。它是几何整射影曲线，函数域为$k(t)$，因而满足原文的不可约、可分等常用解释。原点$x$的局部环是
\[
R=k[t^3,t^4,t^5]_{(t^3,t^4,t^5)}.
\]
其极大理想$\mathfrak m$由三个所列单项式生成；$\mathfrak m^2$中所有非零单项式的指数至少为6，因此三个生成元在$\mathfrak m/\mathfrak m^2$中线性无关。故
\[
\dim R=1,\qquad\dim_k\mathfrak m/\mathfrak m^2=3.
\]
取$m=2$，于是$\dim X=1\leq m-1$。泛中心避开$X$，投影定义在整个$X_{k(\eta)}$上；其概形论像$Y_\eta$是$\mathbb P^2_{k(\eta)}$的闭子概形，在$p(x)$处的切空间维数至多2。若$p'_{C_\eta}$在$x_{k(\eta)}$处非分歧，几何切映射就应单射；但源维数3、目标至多2，矛盾。因此它更不可能在该点平展。

这个点位于原来的$k$有理闭点$x$上方，完全属于原稿的量词范围。反例不借助非约化源、不借助不可分函数域，也不借助可约性。原14.5的双有理部分与这个局部平展断言是不同结论；否定后者并不否定前者。

\audit{V22-DIF-E05}{14.5的订正范围与独立补证}{保留双有理结论；所涉闭点光滑时有局部同构}
本项取有限型的几何约化$k$-概形$X$，浸入$X\hookrightarrow\mathbb P^r_k$，$\dim X\leq m-1$。这是本版对原“可分”的明确工作范围，不将该术语擅自解释为处处光滑。设$\overline X$为约化射影闭包。$X$在$\overline X$中为开子概形；每个分支的函数域可分，故$\overline X$仍几何约化。有限型几何约化概形有光滑稠密开集，参见\href{https://stacks.math.columbia.edu/tag/056V}{Stacks Project，Tag 056V}。

\lead{第一步：有限投影。}
泛中心避开$\overline X$，因为$\dim\overline X\leq m-1$。线性投影的每条纤维是某个包含中心的线性空间去掉其中的超平面，因而是仿射空间。它与$\overline X$的交既为固有概形又为仿射概形，所以有限。投影$\overline X_{k(\eta)}\to\mathbb P^m_{k(\eta)}$固有且拟有限，故有限。此处同时用了中心避开闭包与闭包的固有性；不能只由$X$拟有限就断言$X$的投影有限。

\lead{第二步：一个原光滑闭点上方没有第二个几何点。}
固定$X$的一个光滑闭点$t$，先作代数闭扩域，处理其有限多个几何点。对其中一点$a$，令$b$遍历$\overline X\setminus\{a\}$。把$a,b$投到同一点的中心满足与割线$\overline{ab}$相交的条件。上述关联维数至多
\[
N+\dim\overline X-m\leq N-1.
\]
因此泛中心避开所有这样的割线坏条件，也分离$t$的不同几何共轭。再排除$a$处的切方向坏条件，得到微分单射。$X$与$\overline X$在$t$附近相同，故这里光滑切空间维数至多$m-1$。有限个几何点的好条件按E02、E03下降，不要求$t$在$k$上有理。

\lead{第三步：有限性、剩余域与Nakayama步骤。}
在$y=p'(x)$的纤维上，第二步给出唯一几何点，因而剩余域扩张纯不可分；非分歧性又使它有限可分，所以剩余域相同，且纤维为一个约化点。先从目标删去非分歧坏集的有限像，投影就在$y$附近处处非分歧。有限环映射$A\to B$局部满足
\[
B/\mathfrak m_A B=\kappa(y).
\]
对有限$A$-模$B/\operatorname{im}A$应用Nakayama引理，得到局部满射$A\to B$。故投影在$y$的邻域是闭浸入；这是\href{https://stacks.math.columbia.edu/tag/04DG}{Stacks Project，Tag 04DG}的局部步骤。目标取概形论像时，该环映射还单射，因此成为同构。由于整个闭包纤维只有属于$X$的点，亦可删去闭包边界的有限像，使此邻域完全来自$X$。所以原光滑闭点上方的投影实际上在其像的邻域给出同构，特别地平展；其整条逆像也包含在该同构邻域中。

这里没有用“几何单点纤维”单独推出双有理或平展。纯不可分的有限态射也可能有单点纤维；微分及剩余域步骤不可删去。对每一个固定的原光滑闭点$t$，上述过程给出一个非空好开集，泛点$\eta$自然属于它。因此结论可以逐个用于所有原光滑闭点；无需声称无穷多个好开集的交仍为开集或仍有$k$点，也不能把量词扩大为$X_{k(\eta)}$的全部闭点。

\lead{第四步：各分支的双有理性。}
在$\overline X_{\bar k}$的每个不可约分支上，取一个位于$X$的光滑稠密开内、又不属于其他分支的几何点。分支和点的数目有限，这些选择定义在某个有限扩域上。对全部点同时应用前述好中心条件，得到非空开集。在每一点的像附近，投影是同构；没有别的分支映到这个局部像中。这保证各分支的像互异，且在每个分支上存在非空开集同构，故相应函数域相同。结论沿域的忠实平坦扩张下降，得到泛中心投影的双有理性。这个做法也适用于分支维数不同的情形，不以总空间的最大维数计数来掩盖较小分支。

\audit{V22-DIF-E06}{14.6与由泛中心到一般中心}{给出明确非空开条件，不把被否定的局部断言用作前提}
令$X$为原14.6中的不可约、可分射影代数概形，按几何约化解释，则$X$整。必要时先提高射影环境维数，使$r\geq n+1$，取$m=n+1$。上项各个光滑指定点的割线回避、切方向回避和中心避开$X$的条件都是非空开条件；在某有限扩域定义后，可删去其坏闭集在$k$上Grassmann概形中的有限像。这些像维数仍小于$N$，故得到定义在$k$上的非空开子概形$U\subset\mathcal G$，每个几何中心都给出有限双有理投影。

Grassmann概形有一个同构于仿射空间的稠密大胞腔。当$k$无限时，仿射空间的$k$有理点Zariski稠密，故$U(k)\ne\varnothing$；若中心为空，则唯一中心直接给出结论。选择这样的$k$中心$C$，得到
\[
X\longrightarrow Y\subset\mathbb P^{n+1}_k
\]
有限双有理且满射。$Y$是$n$维整闭子概形，在射影空间中余维1。它的齐次素理想是多项式唯一分解环中的高度1齐次素理想，因此由一个齐次不可约多项式生成；$Y$是超曲面。零维时同样适用，可分整射影零维概形是有限可分域扩张，这个论证给出$\mathbb P^1$中的有限超曲面。

这只证明原来在无限域上的结论；没有断言任意非空开集在有限域上有有理点。也没有将“可构造”误作“每个点都成立”：若不可约诺特空间$G$的可构造子集$A$含泛点，将$A$写成有限个局部闭集之并，含泛点的那项的闭包必为$G$，于是其中含一个非空开集。只有当所讨论的整条纤维性质确实给出这样的参数可构造集时，原稿的邻近纤维推论才成立。

\audit{V22-DIF-E07}{作者关于非浸入、非非分歧的警告}{保留原语气并给出可核算的曲线补例}
在$\mathbb C$上，曲线$y^2=x^5-x$的光滑射影模型是亏格2曲线：到$\mathbb P^1$的二重覆盖在五个互异有限根及无穷点共六点分歧，Riemann--Hurwitz公式给出$2g-2=-4+6=2$。若它浸入$\mathbb P^2$，固有浸入必为闭浸入，像是光滑平面曲线；次数$d$的光滑平面曲线亏格为$(d-1)(d-2)/2$，没有整数$d$给出2。因此它不能嵌入$\mathbb P^2$。这具体说明14.6的双有理态射不能一般加强成浸入。

关于到$\mathbb P^m$的非分歧性，也可在$m=1$取任意光滑射影椭圆曲线。有限满射$E\to\mathbb P^1$若在特征0处处非分歧，就会是有限平展；Riemann--Hurwitz给出$0=-2\deg(f)$，矛盾。这个补例已经独立否定无条件的全局非分歧断言，不必伪称本卷已证明所引“第VIII章”的全部基本群结果。

原文的$X\to\mathbb P^m$同构论证仍保留固有、连通、光滑等维和代数单连通的用途。正文14.4控制的是指定的原光滑点；它不声称泛扩域后的每个闭点都来自一个原闭点，也不把一般投影的分歧集判成空集。

\audit{V22-DIF-E08}{保存与未决边界}{旧文逐字保留；本轮没有替作者完成全部微分奇点理论}
旧后半段及V-H-L029--047共19条旧注逐字保存为本轮历史副本，原continuations/v\_dot\_late.tex整体仍在继承源码中。新正文连续翻译法88第二段至法90末本节结束；原14.4(b)的手插加强和14.5的原量词都在新脚注中明确保留并订正，未静默删去错误。

尚未完成的是微分性质导言所援引的第12节整组普通奇点、唯一性及有理性定理的全一般性审证，以及原“可分”在所有历史上下文中的统一术语追溯；本轮充分证明明确使用有限型和几何约化。下一回校点为法90末、英59的第15节。第15节旧译及后续篇章全部保留，本轮没有将它们算作已经重新回校。
